\documentclass[11pt]{article}
\usepackage{mathblatt}
% Präambel des Setzers (werkzeuge/setzer-praeambel.tex, wortgleich); dazu \gteil (Muster T2U)
\makeatletter
\setlength{\mbpfbe}{0mm}% keine Punkte (pruefung.md § 5: nur beim Prüfstein)
\newlength{\gEin}\setlength{\gEin}{0pt}
\let\gSchrift\relax
\renewcommand{\pfkreuzzeile}[1]{\par\noindent\hangindent3.2em\hangafter1\hspace*{1.6em}$\square$\ #1\par}
\newcommand{\gkreuz}[1]{$\square$\ #1\hspace{2em}}
\newcommand{\gkreuzzeile}[2]{\par\noindent\hangindent2.2em\hangafter1\mbox{#1\,$\square$}\ #2\par}% Buchstabe, Option: eine je Zeile (neben Skizze)
% links, was man ansieht, rechts, was man tut; Spaltengrenze überall an derselben Stelle (gemeinsam.md „Satz“)
% v0.17: Fußeintrag „Nr: Wert“ – kurz in einer Box (nie getrennt), lang umbrechbar mit der Nummer am Anfang
% (vorher \mbox: lange Lösungen liefen über den Rand, AS40 138 pt); robust, weil \fusshilfe in die .aux schreibt
\NewDocumentCommand{\gfuss}{m m}{\sbox0{#1:~#2}\ifdim\wd0<0.45\textwidth\mbox{#1:~#2}\else\mbox{#1:}\nobreak\ #2\fi}
\newsavebox{\gbox}\newlength{\gLinks}
% v0.10 (Lehrer 13.10.): keine starre Halbierung – die linke Spalte ist so breit wie die Figur (+ 6 mm), die
% Antwort steht direkt daneben; ist die Figur breiter als 0,62 Zeile, steht die Antwort darunter
\newcommand{\gzwei}[2]{\par\smallskip\noindent\sbox{\gbox}{#1}%
  \setlength{\gLinks}{\dimexpr\wd\gbox+6mm\relax}%
  \ifdim\gLinks>0.62\linewidth
    \usebox{\gbox}\par\smallskip\noindent\begin{minipage}{\linewidth}#2\end{minipage}%
  \else
    \begin{minipage}[t]{\gLinks}\vspace{0pt}\usebox{\gbox}\end{minipage}%
    \begin{minipage}[t]{\dimexpr\linewidth-\gLinks\relax}\vspace{0pt}\raggedright #2\end{minipage}%
  \fi\par}
% v0.9 (12.10., pruefung.md § 5 „Marke mit B, Bedingungen rechts“): drittes Argument von pfaufg = Bedingungen
% der Prüfung (v0.10: \pfohnetr = TR durchgestrichen, „FOR“), klein grau im rechten Rand auf der Grundlinie der Nummer, Zeilen mit \\
\newcommand{\gbedrand}[2]{\ifblank{#2}{}{\rlap{\hspace*{\dimexpr#1+2mm\relax}%
  \parbox[t]{15mm}{\raggedright\scriptsize\leavevmode\color{mbgrau}\strut#2}}}}
% Teil eines Bündels (a, b …): Marke links im Rand, Bedingungen rechts im Rand, auf der Zeile des Teils
\newcommand{\gteilrand}[2]{\llap{\mbpfrandmarke{#1}\hspace*{\mbpfnr}}\gbedrand{\linewidth}{#2}}
% Teil eines Bündels: Buchstabe in eigener Spalte, Inhalt eingerückt (\gEin = Einzug, damit die Spaltengrenze von
% \gzwei an derselben Stelle bleibt)
\newenvironment{gbteil}[3]{\par\smallskip\noindent\gteilrand{#2}{#3}\makebox[\mbpfteil][l]{\textbf{#1}}%
  \begin{minipage}[t]{\dimexpr\linewidth-\mbpfteil\relax}\setlength{\gEin}{\mbpfteil}\raggedright\strut\ignorespaces}%
 {\end{minipage}\par}
% Grafik höchstens so breit wie die Spalte (Bündel nebeneinander, Koordinatensystem ohne Karo)
\newsavebox{\gpbox}
\newcommand{\gpassend}[1]{\sbox{\gpbox}{#1}\ifdim\wd\gpbox>\linewidth\typeout{GPASSEND: verkleinert \the\wd\gpbox\space auf \the\linewidth}\resizebox{\linewidth}{!}{\usebox{\gpbox}}\else\usebox{\gpbox}\fi}
% Skizzenfeld (v0.10, Lehrer 13.10.): Karo ab dem linken Rand über die Zeile, „Skizze“ klein grau oben links im Feld
% (wie \pfskizzenfeld der Vorlage); der rechte Rand bleibt den Bedingungen
\newcommand{\gskizze}[1]{\par\smallskip\noindent\begin{tikzpicture}
  \clip (0,0) rectangle ({\linewidth-1mm},{#1*10mm});
  \draw[black!16,line width=0.3pt,step=5mm] (0,0) grid ({\linewidth},{#1*10mm});
  \node[anchor=north west,inner sep=2pt,font=\scriptsize,text=mbgrau,fill=white] at (0,{#1*10mm}) {Skizze};
  \end{tikzpicture}\par}
\RenewDocumentEnvironment{pfaufg}{m m m}{%
  \begin{lrbox}{\mbaufgabebox}\begin{minipage}[t]{\linewidth}%
  \gSchrift\noindent\hspace*{\gEin}\mbpfrandmarke{#1}\gbedrand{\linewidth-\gEin}{#3}%
  \begin{minipage}[t]{\mbpfnr}\strut\textbf{#2}\end{minipage}%
  \begin{minipage}[t]{\dimexpr\linewidth-\gEin-\mbpfnr\relax}\raggedright\strut\ignorespaces}%
 {\end{minipage}%
  \end{minipage}\end{lrbox}%
  \setlength{\mbaufgabehoehe}{\dimexpr\ht\mbaufgabebox+\dp\mbaufgabebox\relax}%
  \ifdefempty{\gkopfnach}{\par\addvspace{5pt}\Needspace*{\mbaufgabehoehe}}{%
    % Stufenkopf, der mit dieser (ersten) Aufgabe zusammenbleibt: Platz für beide verlangen, dann den Kopf setzen
    % (12.10., Quadratische: ein Kopf stand allein unten, die Aufgabe mit Koordinatensystem auf der nächsten Seite)
    % v0.17 (AS40 „Terme deuten“ allein unten): hinter dem Kopf steht das Sprungziel (\special, ein Knoten, nach
    % dem die Seite brechen durfte) – darum \nobreak nach dem Kopf; Vorhalt 56 pt = Kopf mit Linie und Abständen
    \par\Needspace*{\dimexpr\mbaufgabehoehe+56pt\relax}\gkopfnach\par\nobreak\gdef\gkopfnach{}\addvspace{5pt}}%
  \noindent\usebox{\mbaufgabebox}\par\addvspace{5pt}}
% Stufenkopf vormerken: \gkopfvor{<Kopf>} – die nächste pfaufg setzt ihn, sobald beide zusammen Platz haben
\def\gkopfnach{}
\newcommand{\gkopfvor}[1]{\gdef\gkopfnach{#1}}
% Teilaufgabe mit Skizze links in gzwei: Buchstabe in eigener Spalte, darunter Text und Skizze
\newcommand{\gteil}[2]{\makebox[\mbpfteil][l]{\textbf{#1}}\parbox[t]{\dimexpr0.5\textwidth-\gEin-\mbpfnr-5mm-\mbpfteil\relax}{\raggedright #2}}
% Stufenkopf der Vorlage (mathblatt.sty 2026-10-12: Linie, Abstand, 13 pt, \pfjahre) ohne ihren Vorhalt von
% 14 Zeilen: der Setzer fordert vor dem Kopf 0,16 Seitenhöhe (Kopf + erste Aufgabe), mehr nicht
\setlength{\mbpfstufevorhalt}{0pt}
\makeatother
\newcommand{\karo}[1]{\par\smallskip\noindent\begin{tikzpicture}
  \clip (0,0) rectangle ({\linewidth-1mm},{#1*10mm});
  \draw[black!16,line width=0.3pt,step=5mm] (0,0) grid ({\linewidth},{#1*10mm});
  \end{tikzpicture}\par}
\newcommand{\kk}[1]{$\square$\,#1\hspace{0.9em}}
\newcommand{\linie}[1]{\,{\color{black!45}\rule[-1pt]{#1}{0.4pt}}\,}
\newcommand{\pkt}{\textperiodcentered{}}
% v0.14 (14.10.): Stufenkopf mit fünf Punkten (Option = Jahre für die Bestellkarte) oder „zuletzt 19“
\NewDocumentCommand{\gpunkte}{o m}{#2}
\newcommand{\gzuletzt}[1]{zuletzt #1}
% v0.7 (12.10.): Schülerfrage unter dem Stufenkopf; Ankreuzoption als Bild (Buchstabe und Kästchen darunter)
\newcommand{\gfrage}[1]{\par\nobreak\noindent\hspace*{\mbpfnr}{\small\itshape #1}\par\nobreak\addvspace{2pt}}
\newlength{\gbildbreite}\setlength{\gbildbreite}{0.32\linewidth}% Setzer setzt sie je Reihe (vier Bilder: 0,24)
\newcommand{\gbildoption}[3]{\begin{minipage}[t][#3][b]{\gbildbreite}\centering #2\par\smallskip\mbox{#1\,$\square$}\end{minipage}}
% Lineare (12.10.): Bild mit Buchstabe ohne Kästchen (Diagramm zuordnen); Aussage mit „wahr / falsch“ zum Ankreuzen
\newcommand{\gbildname}[3]{\begin{minipage}[t][#3][b]{\gbildbreite}\centering #2\par\smallskip\mbox{\textbf{#1}}\end{minipage}}
\newcommand{\gwahrfalsch}[2]{\par\noindent\hangindent2.2em\hangafter1\makebox[2.2em][l]{(#1)}#2\hfill\mbox{wahr $\square$\quad falsch $\square$}\par}
\makeatletter\setlength{\mbpfbe}{0mm}\makeatother
\begin{document}

\pfheftstil
\pflaufkopf[Abitur GK \pkt{} mehr]{Säulendiagramm und Erwartungswert}{AS80}\pfseitevon
\pfheftkopf{Säulendiagramm und Erwartungswert}{mehr}
\pfstufenzeile{Erwartungswert am Diagramm \pkt{} Erwartungswert berechnen \pkt{} Wahrscheinlichkeit ablesen}
% 2020-be-gk-B4.2d-m1 (Schwere II, rang 11)
\begin{pfaufg}{GK~’20\\B4.2d}{1.}{}
In einem großen Unternehmen ist ein Drittel der Beschäftigten weiblich. Es werden 40 Beschäftigte zufällig ausgewählt; $X$ ist die Anzahl der weiblichen darunter. Begründe ohne Berechnung von Wahrscheinlichkeiten, dass die Wahrscheinlichkeitsverteilung von $X$ für 13 oder 14 ihren größten Wert hat.
\karo{3}
\fusshilfe{\gfuss{1}{größter Wert bei einem Nachbarn des Erwartungswerts, also 13 oder 14}}
\end{pfaufg}
% 2018-be-gk-B3.2b-m1 (Schwere I, rang 12)
\begin{pfaufg}{GK~’18\\B3.2b}{2.}{}
Die Anzahl $X$ der fehlerhaften Geräte unter 240 zufällig ausgewählten Bildschirmen ist binomialverteilt mit $p = 0{,}15$. Gib die Anzahl fehlerhafter Bildschirme an, die mit der größten Wahrscheinlichkeit auftritt.
\karo{3}
\fusshilfe{\gfuss{2}{größte Wahrscheinlichkeit bei 36}}
\end{pfaufg}
% 2019-be-gk-B4.2c-m1 (Schwere I, rang 13)
\begin{pfaufg}{GK~’19\\B4.2c}{3.}{}
Zwei Geschwister würfeln jeden Tag einmal; mit der Wahrscheinlichkeit $\tfrac{5}{12}$ bringt der ältere den Müll hinunter. Berechne den Erwartungswert dafür, wie oft der ältere in 24 Tagen den Müll hinunterbringt.
\karo{2}
\fusshilfe{\gfuss{3}{10}}
\end{pfaufg}
% 2017-iqb-ga-A11a-m1 (Schwere I, rang 14)
\begin{pfaufg}{IQB~’17\\A11a}{4.}{}
Die Abbildung zeigt die Wahrscheinlichkeitsverteilung einer binomialverteilten Zufallsgröße $X$. Ermittle mit der Abbildung einen Näherungswert für $P(4 \le X \le 6)$.
\par\smallskip\noindent\gpassend{\saeulenab[ymax=0.25,ystep=0.05,yfein=0.01,ylabel={$P(X=k)$},hoehe=2.8,breite=0.5]{0/0.0004,1/0.0042,2/0.0190,3/0.0547,4/0.1104,5/0.1664,6/0.1941,7/0.1792,8/0.1327,9/0.0794,10/0.0385,11/0.0151,12/0.0047,13/0.0012,14/0.0002}}\par
\karo{2}
\fusshilfe{\gfuss{4}{0,47}}
\end{pfaufg}
% 2021-iqb-ga-A2b-m1 (Schwere III, rang 15)
\begin{pfaufg}{IQB~’21\\A2b}{5.}{}
Die Abbildung zeigt die Wahrscheinlichkeitsverteilung einer binomialverteilten Zufallsgröße $Y$ mit $p = 0{,}5$; sie ist symmetrisch zu $k = 11{,}5$. Es gilt $P(Y \le 13) \approx 0{,}80$ und $P(Y = 10) \approx 0{,}14$. Berechne mit diesen Werten einen Näherungswert für $P(Y = 12)$.
\par\smallskip\noindent\gpassend{\saeulenab[ymax=0.2,ystep=0.05,yfein=0.05,ylabel={$P(X=k)$},hoehe=2.8,breite=0.5,ohnezahlen]{4/0.0011,5/0.0040,6/0.0120,7/0.0292,8/0.0584,9/0.0974,10/0.1364,11/0.1612,12/0.1612,13/0.1364,14/0.0974,15/0.0584,16/0.0292,17/0.0120,18/0.0040,19/0.0011}}\par
\karo{4}
\fusshilfe{\gfuss{5}{0,16}}
\end{pfaufg}
% 2019-iqb-ga-A12b-m1 (Schwere II, rang 16)
\begin{pfaufg}{IQB~’19\\A12b}{6.}{}
In einer Urne liegen 4 rote und 6 weiße Kugeln. Es wird zehnmal mit Zurücklegen gezogen; $X$ ist die Anzahl der weißen Kugeln. Begründe ohne Berechnung von Wahrscheinlichkeiten, dass keines der Diagramme A und B die Verteilung von $X$ zeigt.
\par\smallskip\noindent\gpassend{\begin{tabular}[b]{@{}c@{}}\saeulenab[ymax=0.3,ystep=0.1,yfein=0.05,ylabel={$P(X=k)$},hoehe=2.8,breite=0.5]{0/0.0060,1/0.0403,2/0.1209,3/0.2150,4/0.2508,5/0.2007,6/0.1115,7/0.0425,8/0.0106,9/0.0016,10/0.0001}\\ \textsf{A}\end{tabular}\quad\begin{tabular}[b]{@{}c@{}}\saeulenab[ymax=0.6,ystep=0.1,yfein=0.05,ylabel={$P(X=k)$},hoehe=2.8,breite=0.5]{0/0.0002,1/0.0031,2/0.0212,3/0.0849,4/0.2230,5/0.4013,6/0.5016,7/0.4300,8/0.2419,9/0.0806,10/0.0121}\\ \textsf{B}\end{tabular}}\par
\karo{4}
\fusshilfe{\gfuss{6}{6; A: höchste Säule bei 4 statt bei 6; B: Säulen summieren sich auf etwa 2 > 1}}
\end{pfaufg}
% 2020-iqb-ga-A2a-m1 (Schwere II, rang 17)
\begin{pfaufg}{IQB~’20\\A2a}{7.}{}
Die Zufallsgröße $X$ ist binomialverteilt mit $n = 6$ und dem Erwartungswert 2. Berechne $P(X = 6)$.
\karo{4}
\fusshilfe{\gfuss{7}{$\tfrac{1}{729}$}}
\end{pfaufg}
\par\vfill\noindent{\scriptsize\color{mbgrau}Originale: Abi GK 2020 \pkt{} B4.2d, Abi GK 2018 \pkt{} B3.2b, Abi GK 2019 \pkt{} B4.2c, IQB GK 2017 \pkt{} A11a, IQB GK 2021 \pkt{} A2b, IQB GK 2019 \pkt{} A12b, IQB GK 2020 \pkt{} A2a\par}
\end{document}